\scanpage{237}

\deleted{On trouve donc
\[
 H^i(C_{\mathrm{DR}}^*(X,\widehat A_*^*))=
 \begin{cases}H^0(X;S)&\text{si }i=0\\H^i(X,S^+)&\text{si }i\ne0\end{cases}
\]}

\[
\deleted{\unsure{E|\&}}\quad
\begin{cases}
 \operatorname{Fil}^p(\widehat A_*^0)\ \text{flasque pour tt }p\geq1\\
 \pi_i(\widehat A_*^0)=\begin{cases}0&\text{si }i\ne0\\ k&\text{si }i=0\end{cases}
 \quad\text{si }p\leq0
\end{cases}
\]

Considérons la filtration d'algèbre \inserted{½simpl.} sur $\widehat A_*^*$ prolongeant celle de $\widehat A_*^0$ et donnant \deleted{\unsure{aux}} à $\widehat A_*^1$ la filtration 1, donc

\[
 \operatorname{Fil}^p A_*^{*}=\sum_{0\leq i\leq p}\operatorname{Fil}^{p-i}(\widehat A_*^0)\otimes_S\bigwedge^i\deleted{\unsure{\underline\Phi}} \Psi_*
\]

\uline{C'est stable par l'opérateur différentiel}. On a \deleted{\unsure{donc}}

\[
\begin{aligned}
\pi_\ell(\operatorname{Fil}^p(A_*^i))
 &=\deleted{\sum_{0\leq\unsure{j|i}\leq p}}\pi_\ell(\operatorname{Fil}^{p-i}(\widehat A_*^0)\otimes_S\bigwedge^i\Psi_*)\\
 &\hspace{2em}\deleted{\pi_{\ell-i}(\operatorname{Fil}^{p-i}(\widehat A_*^0))}\\
 &=\pi_{\ell-i}(\operatorname{Fil}^{p-i}(\widehat A_*^0))\\
 &\quad\deleted{0\ \text{\unsure{sauf} si }\ell-i=\unsure{p-i}}\quad\deleted{\text{\unsure{i.e.} }\ell=p}\\
 &=\begin{cases}0&\text{si }\deleted{\ell\ne p}\ \text{ou }\unsure{i<p}\ \deleted{\unsure{\ell=p|\ell\ne p}}\deleted{\unsure{p\geq p}}\\[-2pt]
 &\inserted{p\leq i,\ \text{et }\ell=i}\ \deleted{\unsure{\text{et }p\leq}}\\
 k&\text{si }\deleted{\ell=p,\ i=p}\ p\leq i=\ell\end{cases}
\end{aligned}
\]

Donc \deleted{on en conclut la suite spectrale (pour $p$ fixé)} si on regarde $\operatorname{Fil}^p(A_*^*)$ comme \deleted{$E_2^{rs}=H^r(\unsure{\mathcal H}^s(X,\unsure{\mathcal F}$} une résolution ½simpliciale de $\deleted{\unsure{A}}S_*^{[p]}=\operatorname{Fil}^p S$, on voit que tous les termes $\operatorname{Fil}^p(A_*^i)$ de cette résolution \inserted{en degré $<p$} sont flasques, \deleted{\unsure{sauf}} après on a \deleted{\unsure{puis}}\deleted{\unsure{le premier}} $(j\geq p)$ $\operatorname{Fil}^p(A_*^j)\simeq\widehat A_*^0\otimes_S\bigwedge^j\Psi_*$ dont le seul groupe d'homotopie non nul est $\pi_j(\cdot)\simeq k$. On en conclut

\diagram{\begin{tikzcd}[ampersand replacement=\&,column sep=small]
 H^i\bigl(\underbrace{\operatorname{Fil}^p C_{\mathrm{DR}}^*(X,\widehat A_*^*)}_{\substack{\text{\unsure{2\textsuperscript{e} défn|cf. défn|$\simeq$ défn}}\\(C_{\mathrm{DR}}^*(X,\operatorname{Fil}\widehat A_*^*))}}\bigr)
 \arrow[r,"\sim"]\arrow[d] \& H^i(X,\underbrace{S^{[p]}}_{=\operatorname{Fil}^p S})\arrow[d]\quad\text{si }i\leq p\\
 H^i_{\unsure{1}}(C_{\mathrm{DR}}^*(X,\widehat A_*^*)\arrow[r] \& H^i(X,S)
\end{tikzcd}}

% archive sheet number: 236
